% GENERATED FILE. Edit canonical lessons, then run npm run generate:books.
% canonical-source-hash: fnv1a64:ee4e2f75c89518d5
% canonical-lessons: UR-C252-papita, UR-C252-khubani, UR-C252-anannas, UR-C252-nariyal, UR-C252-kaju

\chapter{Papaya, Apricot, Pineapple, Coconut, Cashew}
\label{ch:ur-papaya-apricot-pineapple-coconut-cashew}
\hlchaptermodality{252}

\begin{chapteropening}
\textbf{By the end of this chapter:} \emph{I can say a papaya, an apricot, a pineapple, a coconut and a cashew.}

Complete the last lesson of chapter 252: I can say a papaya, an apricot, a pineapple, a coconut and a cashew.
\end{chapteropening}

\section[\texorpdfstring{\ur{پپیتا} (papītā)}{پپیتا (papītā)}]{\ur{پپیتا} (papītā) — a papaya}

\label{lesson:UR-C252-papita}



\begin{quote}
Before the new one: say the Urdu for shorts, then the Urdu for a raincoat.
\end{quote}

\subsection*{You'll want to know: \ur{پپیتا}}
\textbf{\ur{پپیتا}} — \emph{papītā} — \textquotedblleft{}a papaya\textquotedblright{}.

\begin{cousinweb}
The origin of \textbf{\ur{پپیتا}} (\emph{papītā}) is uncertain. It is usually linked with \emph{papaya}, a Caribbean word that Spanish and Portuguese sailors carried round the world; CDIAL has no entry for it.
\end{cousinweb}

\begin{grammarlens}[title={What you've built}]
One more word for this chapter.
\end{grammarlens}

\subsection*{Your turn}
\begin{itemize}
\raggedright
  \item \emph{Say it:} \emph{papītā}
  \item \emph{Say it:} \emph{papītā}, once more
  \item \emph{Say it:} say \emph{papītā}
  \item \emph{Recall:} say the Urdu for shorts, then the Urdu for a raincoat, then say \emph{papītā} again
\end{itemize}

\begin{tcolorbox}[breakable,colback=teal!4,colframe=teal!35!black,title={Before you move on}]
What does \ur{پپیتا} mean? (\textquotedblleft{}A papaya\textquotedblright{}.) Say it once more.
\end{tcolorbox}
\section[\texorpdfstring{\ur{خوبانی} (khūbānī)}{خوبانی (khūbānī)}]{\ur{خوبانی} (khūbānī) — an apricot}

\label{lesson:UR-C252-khubani}



\begin{quote}
Before the new one: say the Urdu for a raincoat, then the Urdu for a papaya.
\end{quote}

\subsection*{You'll want to know: \ur{خوبانی}}
\textbf{\ur{خوبانی}} — \emph{khūbānī} — \textquotedblleft{}an apricot\textquotedblright{}.

\begin{cousinweb}
\textbf{\ur{خوبانی}} (\emph{khūbānī}) is a loan, usually traced to Persian \emph{khūbānī}, \textquotedblleft{}apricot\textquotedblright{}.
\end{cousinweb}

\begin{grammarlens}[title={What you've built}]
One more word for this chapter.
\end{grammarlens}

\subsection*{Your turn}
\begin{itemize}
\raggedright
  \item \emph{Say it:} \emph{khūbānī}
  \item \emph{Say it:} \emph{khūbānī}, once more
  \item \emph{Say it:} say \emph{khūbānī}
  \item \emph{Recall:} say the Urdu for a raincoat, then the Urdu for a papaya, then say \emph{khūbānī} again
\end{itemize}

\begin{tcolorbox}[breakable,colback=teal!4,colframe=teal!35!black,title={Before you move on}]
What does \ur{خوبانی} mean? (\textquotedblleft{}An apricot\textquotedblright{}.) Say it once more.
\end{tcolorbox}
\section[\texorpdfstring{\ur{اناناس} (anannās)}{اناناس (anannās)}]{\ur{اناناس} (anannās) — a pineapple}

\label{lesson:UR-C252-anannas}



\begin{quote}
Before the new one: say the Urdu for a papaya, then the Urdu for an apricot.
\end{quote}

\subsection*{You'll want to know: \ur{اناناس}}
\textbf{\ur{اناناس}} — \emph{anannās} — \textquotedblleft{}a pineapple\textquotedblright{}.

\begin{cousinweb}
\textbf{\ur{اناناس}} (\emph{anannās}) is usually traced to Portuguese \emph{ananás}, \textquotedblleft{}pineapple\textquotedblright{}, which the Portuguese took from a South American (Tupi) word.
\end{cousinweb}

\begin{grammarlens}[title={What you've built}]
One more word for this chapter.
\end{grammarlens}

\subsection*{Your turn}
\begin{itemize}
\raggedright
  \item \emph{Say it:} \emph{anannās}
  \item \emph{Say it:} \emph{anannās}, once more
  \item \emph{Say it:} say \emph{anannās}
  \item \emph{Recall:} say the Urdu for a papaya, then the Urdu for an apricot, then say \emph{anannās} again
\end{itemize}

\begin{tcolorbox}[breakable,colback=teal!4,colframe=teal!35!black,title={Before you move on}]
What does \ur{اناناس} mean? (\textquotedblleft{}A pineapple\textquotedblright{}.) Say it once more.
\end{tcolorbox}
\section[\texorpdfstring{\ur{ناریل} (nāriyal)}{ناریل (nāriyal)}]{\ur{ناریل} (nāriyal) — a coconut}

\label{lesson:UR-C252-nariyal}



\begin{quote}
Before the new one: say the Urdu for an apricot, then the Urdu for a pineapple.
\end{quote}

\subsection*{You'll want to know: \ur{ناریل}}
\textbf{\ur{ناریل}} — \emph{nāriyal} — \textquotedblleft{}a coconut\textquotedblright{}.

\begin{cousinweb}
\textbf{\ur{ناریل}} (\emph{nāriyal}) is a tadbhava, a word inherited from Sanskrit and reshaped by sound change. It comes from \emph{nārikela}, \textquotedblleft{}coconut\textquotedblright{} (CDIAL 7075). CDIAL takes the Sanskrit word from Dravidian.
\end{cousinweb}

\begin{grammarlens}[title={What you've built}]
One more word for this chapter.
\end{grammarlens}

\subsection*{Your turn}
\begin{itemize}
\raggedright
  \item \emph{Say it:} \emph{nāriyal}
  \item \emph{Say it:} \emph{nāriyal}, once more
  \item \emph{Say it:} say \emph{nāriyal}
  \item \emph{Recall:} say the Urdu for an apricot, then the Urdu for a pineapple, then say \emph{nāriyal} again
\end{itemize}

\begin{tcolorbox}[breakable,colback=teal!4,colframe=teal!35!black,title={Before you move on}]
What does \ur{ناریل} mean? (\textquotedblleft{}A coconut\textquotedblright{}.) Say it once more.
\end{tcolorbox}
\section[\texorpdfstring{\ur{کاجو} (kājū)}{کاجو (kājū)}]{\ur{کاجو} (kājū) — a cashew}

\label{lesson:UR-C252-kaju}



\begin{quote}
Before the new one: say the Urdu for a pineapple, then the Urdu for a coconut.
\end{quote}

\subsection*{You'll want to know: \ur{کاجو}}
\textbf{\ur{کاجو}} — \emph{kājū} — \textquotedblleft{}a cashew\textquotedblright{}.

\begin{cousinweb}
\textbf{\ur{کاجو}} (\emph{kājū}) is usually traced to Portuguese \emph{caju}, \textquotedblleft{}cashew\textquotedblright{}, which the Portuguese took from a South American (Tupi) word.
\end{cousinweb}

\begin{grammarlens}[title={What you've built}]
One more word for this chapter.
\end{grammarlens}

\subsection*{Your turn}
\begin{itemize}
\raggedright
  \item \emph{Say it:} \emph{kājū}
  \item \emph{Say it:} \emph{kājū}, once more
  \item \emph{Say it:} say \emph{kājū}
  \item \emph{Recall:} say the Urdu for a pineapple, then the Urdu for a coconut, then say \emph{kājū} again
\end{itemize}

\begin{tcolorbox}[breakable,colback=teal!4,colframe=teal!35!black,title={Before you move on}]
What does \ur{کاجو} mean? (\textquotedblleft{}A cashew\textquotedblright{}.) Say it once more.
\end{tcolorbox}
